\documentclass[aspectratio=169,10pt]{beamer}
% ---------------------------------------------------------------
% R3.07 SQL dans un langage de programmation -- preambule commun des diapositives
% Charte derivee de R5.A.08 Qualite de developpement (metropolis, dark tech)
% Palette IUT partagee (violet du logo + accent or)
% Utilisation :
%   \documentclass[aspectratio=169,10pt]{beamer}
%   % ---------------------------------------------------------------
% R3.07 SQL dans un langage de programmation -- preambule commun des diapositives
% Charte derivee de R5.A.08 Qualite de developpement (metropolis, dark tech)
% Palette IUT partagee (violet du logo + accent or)
% Utilisation :
%   \documentclass[aspectratio=169,10pt]{beamer}
%   % ---------------------------------------------------------------
% R3.07 SQL dans un langage de programmation -- preambule commun des diapositives
% Charte derivee de R5.A.08 Qualite de developpement (metropolis, dark tech)
% Palette IUT partagee (violet du logo + accent or)
% Utilisation :
%   \documentclass[aspectratio=169,10pt]{beamer}
%   \input{../../latex/common/r307-slides-preamble.tex}
%   \graphicspath{{../../latex/common/assets/}{img/}{figs/}}
% ---------------------------------------------------------------
\usetheme{metropolis}
\usepackage{appendixnumberbeamer}

\usepackage[T1]{fontenc}
\usepackage[utf8]{inputenc}
\usepackage[french]{babel}
\usepackage[scale=0.94]{tgheros}
\renewcommand{\familydefault}{\sfdefault}
\usepackage{amsmath,amssymb,amsfonts}
\usepackage{graphicx}
\usepackage{booktabs}
\usepackage{multirow}
\usepackage{tikz}
\usetikzlibrary{shapes.geometric, shapes, arrows, arrows.meta, positioning, calc, fit}
\usepackage{xcolor}
\usepackage{listings}
\usepackage[skins,listings]{tcolorbox}
\usepackage{qrcode}
\usepackage{hyperref}

% ---------------------------------------------------------------
% PALETTE (dark-mode tech: near-black canvas, violet (logo IUT) + gold accents)
% ---------------------------------------------------------------
\definecolor{qabg}{HTML}{0E1016}      % canvas background
\definecolor{qapanel}{HTML}{181B24}   % card / block panel
\definecolor{qapanelalt}{HTML}{20242F}% secondary panel tint
\definecolor{qaline}{HTML}{2B303D}    % hairline rules / table rules
\definecolor{qatext}{HTML}{ECEEF2}    % primary text (off-white)
\definecolor{qamuted}{HTML}{8B93A6}   % secondary / caption text
\definecolor{qaaccent}{HTML}{D93ED7}  % primary accent -- violet du logo IUT, eclairci pour fond sombre
\definecolor{qaaccent2}{HTML}{F2B84B} % secondary accent -- or/ambre
\definecolor{qacodekw}{HTML}{C792EA}  % code : mots-cles
\definecolor{qacodestr}{HTML}{C3E88D} % code : chaines
\definecolor{qacodecom}{HTML}{7F8AA3} % code : commentaires
\definecolor{qacodeann}{HTML}{F2B84B} % code : annotations / javadoc

\setbeamercolor{background canvas}{bg=qabg}
\setbeamercolor{normal text}{fg=qatext}
\setbeamercolor{palette primary}{bg=qabg, fg=qamuted}
\setbeamercolor{frametitle}{bg=qabg, fg=qatext}
\setbeamercolor{progress bar}{fg=qaaccent, bg=qaline}
\setbeamercolor{alerted text}{fg=qaaccent2}
\setbeamercolor{block title}{bg=qapanel, fg=qaaccent}
\setbeamercolor{block body}{bg=qapanel, fg=qatext}
\setbeamercolor{block title alerted}{bg=qapanel, fg=qaaccent2}
\setbeamercolor{block body alerted}{bg=qapanel, fg=qatext}
\setbeamercolor{block title example}{bg=qapanelalt, fg=qaaccent}
\setbeamercolor{block body example}{bg=qapanelalt, fg=qatext}
\setbeamercolor{example text}{fg=qaaccent}
\setbeamercolor{itemize item}{fg=qaaccent}
\setbeamercolor{itemize subitem}{fg=qaaccent2}
\setbeamercolor{enumerate item}{fg=qaaccent}
\setbeamercolor{enumerate subitem}{fg=qaaccent2}
\setbeamercolor{caption}{fg=qamuted}
\setbeamercolor{page number in head/foot}{fg=qamuted, bg=qabg}
\setbeamercolor{title}{fg=qatext}
\setbeamercolor{subtitle}{fg=qamuted}
\setbeamercolor{author}{fg=qatext}
\setbeamercolor{institute}{fg=qamuted}
\setbeamercolor{date}{fg=qamuted}
\setbeamercolor{section in toc}{fg=qatext}
\setbeamercolor{subsection in toc}{fg=qamuted}

\setbeamertemplate{blocks}[rounded][shadow=false]
\setbeamertemplate{itemize item}{\raisebox{0.15em}{\scriptsize\textcolor{qaaccent}{\ensuremath{\blacktriangleright}}}}
\setbeamertemplate{itemize subitem}{\raisebox{0.1em}{\tiny\textcolor{qaaccent2}{\ensuremath{\bullet}}}}

\setbeamerfont{frametitle}{series=\bfseries, size=\Large}
\setbeamerfont{title}{series=\bfseries, size=\Large}
\setbeamerfont{subtitle}{size=\normalsize}
\setbeamerfont{institute}{size=\scriptsize}
\setbeamerfont{date}{size=\scriptsize}
\setbeamerfont{author}{size=\small}

\metroset{sectionpage=none, subsectionpage=none, numbering=fraction, progressbar=frametitle}
\setbeamertemplate{section in toc}[sections numbered]

% ---------------------------------------------------------------
% LISTINGS -- theme sombre (lstlisting reste utilisable tel quel)
% \lstcode : police + couleur de base, a reutiliser dans les
%            surcharges locales (ex. basicstyle=\lstcode\scriptsize)
% ---------------------------------------------------------------
\newcommand{\lstcode}{\ttfamily\color{qatext}}

\lstdefinelanguage{TypeScript}{
  keywords={typeof, new, true, false, catch, function, return, null, switch, var,
    if, in, while, do, else, case, break, const, let, interface, class, extends,
    implements, enum, export, import, this, super, constructor, public, private,
    protected, static, async, await, throw, throws, try, finally, void},
  ndkeywords={boolean, any, number, string, int, long, double},
  sensitive=true,
  comment=[l]{//},
  morecomment=[s]{/*}{*/},
  morestring=[b]',
  morestring=[b]"
}

\lstset{
  language=Java,
  basicstyle=\lstcode\footnotesize,
  keywordstyle=\color{qacodekw}\bfseries,
  ndkeywordstyle=\color{qaaccent2}\bfseries,
  identifierstyle=\color{qatext},
  stringstyle=\color{qacodestr},
  commentstyle=\color{qacodecom}\itshape,
  morecomment=[s][\color{qacodeann}]{/**}{*/},
  backgroundcolor=\color{qapanelalt},
  rulecolor=\color{qaline},
  frame=single,
  framesep=3pt,
  xleftmargin=4pt,
  xrightmargin=4pt,
  keepspaces=true,
  columns=fullflexible,
  showstringspaces=false,
  breaklines=true,
  tabsize=2,
  escapeinside={(*@}{@*)},
}

% ---------------------------------------------------------------
% CARD-STYLE BOXES
% ---------------------------------------------------------------
\tcbset{
  figcard/.style={colback=white, colframe=white, boxrule=0pt, arc=4pt,
    left=3pt, right=3pt, top=3pt, bottom=3pt, boxsep=0pt,
    halign=center, valign=center},
  datacard/.style={colback=qapanel, colframe=qaline, boxrule=0.4pt, arc=3pt,
    left=8pt, right=8pt, top=6pt, bottom=6pt, coltext=qatext,
    before upper={\color{qatext}}, halign=center},
  codebox/.style={colback=qapanelalt, colframe=qaline, boxrule=0.4pt, arc=3pt,
    left=8pt, right=8pt, top=5pt, bottom=5pt, listing only,
    listing options={basicstyle=\ttfamily\small\color{qaaccent},
      breaklines=true, showstringspaces=false, columns=fullflexible,
      keepspaces=true, frame=none, backgroundcolor={}}}
}
\newtcolorbox{figcard}[1][]{figcard,#1}
\newtcolorbox{datacard}{datacard}
\newtcblisting{codebox}{codebox}

% ---------------------------------------------------------------
% HANDY MACROS
% ---------------------------------------------------------------
\newcommand{\hl}[1]{\alert{#1}}
\newcommand{\source}[1]{{\scriptsize\color{qamuted}#1}}

% Sommaire automatique en debut de section
\AtBeginSection[]{%
  \begin{frame}<beamer>{Plan du cours}
    \tableofcontents[currentsection]
  \end{frame}}

% Page de titre : \rtitleslide{numero de chapitre}{titre du chapitre}
% Les auteurs sont repris de \rauthors (modifiable avant \begin{document})
\newcommand{\rauthors}{Selena Lamari et Chouki Tibermacine}
\newcommand{\rtitleslide}[2]{%
\begin{frame}[plain,noframenumbering]
  \begin{center}
    \vspace{2mm}
    {\color{qaaccent}\rule{2.1cm}{1.2pt}}\\[6mm]
    {\Large\bfseries\color{qatext} R3.07}\\[1mm]
    {\Large\bfseries\color{qatext} SQL dans un langage de programmation}\\[5mm]
    {\large\color{qamuted} Chapitre #1}\\[1mm]
    {\large\color{qaaccent} #2}\\[6mm]
    \includegraphics[height=2.1cm]{iut_logo_transparent.png}

    \vspace{4mm}

    {\small\color{qatext} \rauthors}\\[1mm]
    {\scriptsize\color{qamuted} BUT 2 Informatique --- IUT de Vannes --- Universit\'e Bretagne Sud}
  \end{center}
\end{frame}
}

% ---- Diapo de fin « Questions ? » (même modèle que R5.A.08) ----
\newcommand{\rcontact}{Chouki.Tibermacine@univ-ubs.fr}
\newcommand{\rquestionslide}{%
\begin{frame}{Questions ?}
  \begin{center}
    \vspace{2mm}
    \begin{tikzpicture}
      \node[regular polygon, regular polygon sides=6, draw=qaline, thick,
        minimum size=4.6cm, rotate=90] (hexin) {};
      \node[regular polygon, regular polygon sides=6, draw=qaaccent, dashed,
        thick, minimum size=6.2cm, rotate=90] (hexout) {};
      \node[font=\fontsize{92}{92}\selectfont\bfseries, color=qaaccent2] at (hexin) {?};
    \end{tikzpicture}\\[5mm]
    {\small\color{qamuted} \rcontact}
  \end{center}
\end{frame}}

%   \graphicspath{{../../latex/common/assets/}{img/}{figs/}}
% ---------------------------------------------------------------
\usetheme{metropolis}
\usepackage{appendixnumberbeamer}

\usepackage[T1]{fontenc}
\usepackage[utf8]{inputenc}
\usepackage[french]{babel}
\usepackage[scale=0.94]{tgheros}
\renewcommand{\familydefault}{\sfdefault}
\usepackage{amsmath,amssymb,amsfonts}
\usepackage{graphicx}
\usepackage{booktabs}
\usepackage{multirow}
\usepackage{tikz}
\usetikzlibrary{shapes.geometric, shapes, arrows, arrows.meta, positioning, calc, fit}
\usepackage{xcolor}
\usepackage{listings}
\usepackage[skins,listings]{tcolorbox}
\usepackage{qrcode}
\usepackage{hyperref}

% ---------------------------------------------------------------
% PALETTE (dark-mode tech: near-black canvas, violet (logo IUT) + gold accents)
% ---------------------------------------------------------------
\definecolor{qabg}{HTML}{0E1016}      % canvas background
\definecolor{qapanel}{HTML}{181B24}   % card / block panel
\definecolor{qapanelalt}{HTML}{20242F}% secondary panel tint
\definecolor{qaline}{HTML}{2B303D}    % hairline rules / table rules
\definecolor{qatext}{HTML}{ECEEF2}    % primary text (off-white)
\definecolor{qamuted}{HTML}{8B93A6}   % secondary / caption text
\definecolor{qaaccent}{HTML}{D93ED7}  % primary accent -- violet du logo IUT, eclairci pour fond sombre
\definecolor{qaaccent2}{HTML}{F2B84B} % secondary accent -- or/ambre
\definecolor{qacodekw}{HTML}{C792EA}  % code : mots-cles
\definecolor{qacodestr}{HTML}{C3E88D} % code : chaines
\definecolor{qacodecom}{HTML}{7F8AA3} % code : commentaires
\definecolor{qacodeann}{HTML}{F2B84B} % code : annotations / javadoc

\setbeamercolor{background canvas}{bg=qabg}
\setbeamercolor{normal text}{fg=qatext}
\setbeamercolor{palette primary}{bg=qabg, fg=qamuted}
\setbeamercolor{frametitle}{bg=qabg, fg=qatext}
\setbeamercolor{progress bar}{fg=qaaccent, bg=qaline}
\setbeamercolor{alerted text}{fg=qaaccent2}
\setbeamercolor{block title}{bg=qapanel, fg=qaaccent}
\setbeamercolor{block body}{bg=qapanel, fg=qatext}
\setbeamercolor{block title alerted}{bg=qapanel, fg=qaaccent2}
\setbeamercolor{block body alerted}{bg=qapanel, fg=qatext}
\setbeamercolor{block title example}{bg=qapanelalt, fg=qaaccent}
\setbeamercolor{block body example}{bg=qapanelalt, fg=qatext}
\setbeamercolor{example text}{fg=qaaccent}
\setbeamercolor{itemize item}{fg=qaaccent}
\setbeamercolor{itemize subitem}{fg=qaaccent2}
\setbeamercolor{enumerate item}{fg=qaaccent}
\setbeamercolor{enumerate subitem}{fg=qaaccent2}
\setbeamercolor{caption}{fg=qamuted}
\setbeamercolor{page number in head/foot}{fg=qamuted, bg=qabg}
\setbeamercolor{title}{fg=qatext}
\setbeamercolor{subtitle}{fg=qamuted}
\setbeamercolor{author}{fg=qatext}
\setbeamercolor{institute}{fg=qamuted}
\setbeamercolor{date}{fg=qamuted}
\setbeamercolor{section in toc}{fg=qatext}
\setbeamercolor{subsection in toc}{fg=qamuted}

\setbeamertemplate{blocks}[rounded][shadow=false]
\setbeamertemplate{itemize item}{\raisebox{0.15em}{\scriptsize\textcolor{qaaccent}{\ensuremath{\blacktriangleright}}}}
\setbeamertemplate{itemize subitem}{\raisebox{0.1em}{\tiny\textcolor{qaaccent2}{\ensuremath{\bullet}}}}

\setbeamerfont{frametitle}{series=\bfseries, size=\Large}
\setbeamerfont{title}{series=\bfseries, size=\Large}
\setbeamerfont{subtitle}{size=\normalsize}
\setbeamerfont{institute}{size=\scriptsize}
\setbeamerfont{date}{size=\scriptsize}
\setbeamerfont{author}{size=\small}

\metroset{sectionpage=none, subsectionpage=none, numbering=fraction, progressbar=frametitle}
\setbeamertemplate{section in toc}[sections numbered]

% ---------------------------------------------------------------
% LISTINGS -- theme sombre (lstlisting reste utilisable tel quel)
% \lstcode : police + couleur de base, a reutiliser dans les
%            surcharges locales (ex. basicstyle=\lstcode\scriptsize)
% ---------------------------------------------------------------
\newcommand{\lstcode}{\ttfamily\color{qatext}}

\lstdefinelanguage{TypeScript}{
  keywords={typeof, new, true, false, catch, function, return, null, switch, var,
    if, in, while, do, else, case, break, const, let, interface, class, extends,
    implements, enum, export, import, this, super, constructor, public, private,
    protected, static, async, await, throw, throws, try, finally, void},
  ndkeywords={boolean, any, number, string, int, long, double},
  sensitive=true,
  comment=[l]{//},
  morecomment=[s]{/*}{*/},
  morestring=[b]',
  morestring=[b]"
}

\lstset{
  language=Java,
  basicstyle=\lstcode\footnotesize,
  keywordstyle=\color{qacodekw}\bfseries,
  ndkeywordstyle=\color{qaaccent2}\bfseries,
  identifierstyle=\color{qatext},
  stringstyle=\color{qacodestr},
  commentstyle=\color{qacodecom}\itshape,
  morecomment=[s][\color{qacodeann}]{/**}{*/},
  backgroundcolor=\color{qapanelalt},
  rulecolor=\color{qaline},
  frame=single,
  framesep=3pt,
  xleftmargin=4pt,
  xrightmargin=4pt,
  keepspaces=true,
  columns=fullflexible,
  showstringspaces=false,
  breaklines=true,
  tabsize=2,
  escapeinside={(*@}{@*)},
}

% ---------------------------------------------------------------
% CARD-STYLE BOXES
% ---------------------------------------------------------------
\tcbset{
  figcard/.style={colback=white, colframe=white, boxrule=0pt, arc=4pt,
    left=3pt, right=3pt, top=3pt, bottom=3pt, boxsep=0pt,
    halign=center, valign=center},
  datacard/.style={colback=qapanel, colframe=qaline, boxrule=0.4pt, arc=3pt,
    left=8pt, right=8pt, top=6pt, bottom=6pt, coltext=qatext,
    before upper={\color{qatext}}, halign=center},
  codebox/.style={colback=qapanelalt, colframe=qaline, boxrule=0.4pt, arc=3pt,
    left=8pt, right=8pt, top=5pt, bottom=5pt, listing only,
    listing options={basicstyle=\ttfamily\small\color{qaaccent},
      breaklines=true, showstringspaces=false, columns=fullflexible,
      keepspaces=true, frame=none, backgroundcolor={}}}
}
\newtcolorbox{figcard}[1][]{figcard,#1}
\newtcolorbox{datacard}{datacard}
\newtcblisting{codebox}{codebox}

% ---------------------------------------------------------------
% HANDY MACROS
% ---------------------------------------------------------------
\newcommand{\hl}[1]{\alert{#1}}
\newcommand{\source}[1]{{\scriptsize\color{qamuted}#1}}

% Sommaire automatique en debut de section
\AtBeginSection[]{%
  \begin{frame}<beamer>{Plan du cours}
    \tableofcontents[currentsection]
  \end{frame}}

% Page de titre : \rtitleslide{numero de chapitre}{titre du chapitre}
% Les auteurs sont repris de \rauthors (modifiable avant \begin{document})
\newcommand{\rauthors}{Selena Lamari et Chouki Tibermacine}
\newcommand{\rtitleslide}[2]{%
\begin{frame}[plain,noframenumbering]
  \begin{center}
    \vspace{2mm}
    {\color{qaaccent}\rule{2.1cm}{1.2pt}}\\[6mm]
    {\Large\bfseries\color{qatext} R3.07}\\[1mm]
    {\Large\bfseries\color{qatext} SQL dans un langage de programmation}\\[5mm]
    {\large\color{qamuted} Chapitre #1}\\[1mm]
    {\large\color{qaaccent} #2}\\[6mm]
    \includegraphics[height=2.1cm]{iut_logo_transparent.png}

    \vspace{4mm}

    {\small\color{qatext} \rauthors}\\[1mm]
    {\scriptsize\color{qamuted} BUT 2 Informatique --- IUT de Vannes --- Universit\'e Bretagne Sud}
  \end{center}
\end{frame}
}

% ---- Diapo de fin « Questions ? » (même modèle que R5.A.08) ----
\newcommand{\rcontact}{Chouki.Tibermacine@univ-ubs.fr}
\newcommand{\rquestionslide}{%
\begin{frame}{Questions ?}
  \begin{center}
    \vspace{2mm}
    \begin{tikzpicture}
      \node[regular polygon, regular polygon sides=6, draw=qaline, thick,
        minimum size=4.6cm, rotate=90] (hexin) {};
      \node[regular polygon, regular polygon sides=6, draw=qaaccent, dashed,
        thick, minimum size=6.2cm, rotate=90] (hexout) {};
      \node[font=\fontsize{92}{92}\selectfont\bfseries, color=qaaccent2] at (hexin) {?};
    \end{tikzpicture}\\[5mm]
    {\small\color{qamuted} \rcontact}
  \end{center}
\end{frame}}

%   \graphicspath{{../../latex/common/assets/}{img/}{figs/}}
% ---------------------------------------------------------------
\usetheme{metropolis}
\usepackage{appendixnumberbeamer}

\usepackage[T1]{fontenc}
\usepackage[utf8]{inputenc}
\usepackage[french]{babel}
\usepackage[scale=0.94]{tgheros}
\renewcommand{\familydefault}{\sfdefault}
\usepackage{amsmath,amssymb,amsfonts}
\usepackage{graphicx}
\usepackage{booktabs}
\usepackage{multirow}
\usepackage{tikz}
\usetikzlibrary{shapes.geometric, shapes, arrows, arrows.meta, positioning, calc, fit}
\usepackage{xcolor}
\usepackage{listings}
\usepackage[skins,listings]{tcolorbox}
\usepackage{qrcode}
\usepackage{hyperref}

% ---------------------------------------------------------------
% PALETTE (dark-mode tech: near-black canvas, violet (logo IUT) + gold accents)
% ---------------------------------------------------------------
\definecolor{qabg}{HTML}{0E1016}      % canvas background
\definecolor{qapanel}{HTML}{181B24}   % card / block panel
\definecolor{qapanelalt}{HTML}{20242F}% secondary panel tint
\definecolor{qaline}{HTML}{2B303D}    % hairline rules / table rules
\definecolor{qatext}{HTML}{ECEEF2}    % primary text (off-white)
\definecolor{qamuted}{HTML}{8B93A6}   % secondary / caption text
\definecolor{qaaccent}{HTML}{D93ED7}  % primary accent -- violet du logo IUT, eclairci pour fond sombre
\definecolor{qaaccent2}{HTML}{F2B84B} % secondary accent -- or/ambre
\definecolor{qacodekw}{HTML}{C792EA}  % code : mots-cles
\definecolor{qacodestr}{HTML}{C3E88D} % code : chaines
\definecolor{qacodecom}{HTML}{7F8AA3} % code : commentaires
\definecolor{qacodeann}{HTML}{F2B84B} % code : annotations / javadoc

\setbeamercolor{background canvas}{bg=qabg}
\setbeamercolor{normal text}{fg=qatext}
\setbeamercolor{palette primary}{bg=qabg, fg=qamuted}
\setbeamercolor{frametitle}{bg=qabg, fg=qatext}
\setbeamercolor{progress bar}{fg=qaaccent, bg=qaline}
\setbeamercolor{alerted text}{fg=qaaccent2}
\setbeamercolor{block title}{bg=qapanel, fg=qaaccent}
\setbeamercolor{block body}{bg=qapanel, fg=qatext}
\setbeamercolor{block title alerted}{bg=qapanel, fg=qaaccent2}
\setbeamercolor{block body alerted}{bg=qapanel, fg=qatext}
\setbeamercolor{block title example}{bg=qapanelalt, fg=qaaccent}
\setbeamercolor{block body example}{bg=qapanelalt, fg=qatext}
\setbeamercolor{example text}{fg=qaaccent}
\setbeamercolor{itemize item}{fg=qaaccent}
\setbeamercolor{itemize subitem}{fg=qaaccent2}
\setbeamercolor{enumerate item}{fg=qaaccent}
\setbeamercolor{enumerate subitem}{fg=qaaccent2}
\setbeamercolor{caption}{fg=qamuted}
\setbeamercolor{page number in head/foot}{fg=qamuted, bg=qabg}
\setbeamercolor{title}{fg=qatext}
\setbeamercolor{subtitle}{fg=qamuted}
\setbeamercolor{author}{fg=qatext}
\setbeamercolor{institute}{fg=qamuted}
\setbeamercolor{date}{fg=qamuted}
\setbeamercolor{section in toc}{fg=qatext}
\setbeamercolor{subsection in toc}{fg=qamuted}

\setbeamertemplate{blocks}[rounded][shadow=false]
\setbeamertemplate{itemize item}{\raisebox{0.15em}{\scriptsize\textcolor{qaaccent}{\ensuremath{\blacktriangleright}}}}
\setbeamertemplate{itemize subitem}{\raisebox{0.1em}{\tiny\textcolor{qaaccent2}{\ensuremath{\bullet}}}}

\setbeamerfont{frametitle}{series=\bfseries, size=\Large}
\setbeamerfont{title}{series=\bfseries, size=\Large}
\setbeamerfont{subtitle}{size=\normalsize}
\setbeamerfont{institute}{size=\scriptsize}
\setbeamerfont{date}{size=\scriptsize}
\setbeamerfont{author}{size=\small}

\metroset{sectionpage=none, subsectionpage=none, numbering=fraction, progressbar=frametitle}
\setbeamertemplate{section in toc}[sections numbered]

% ---------------------------------------------------------------
% LISTINGS -- theme sombre (lstlisting reste utilisable tel quel)
% \lstcode : police + couleur de base, a reutiliser dans les
%            surcharges locales (ex. basicstyle=\lstcode\scriptsize)
% ---------------------------------------------------------------
\newcommand{\lstcode}{\ttfamily\color{qatext}}

\lstdefinelanguage{TypeScript}{
  keywords={typeof, new, true, false, catch, function, return, null, switch, var,
    if, in, while, do, else, case, break, const, let, interface, class, extends,
    implements, enum, export, import, this, super, constructor, public, private,
    protected, static, async, await, throw, throws, try, finally, void},
  ndkeywords={boolean, any, number, string, int, long, double},
  sensitive=true,
  comment=[l]{//},
  morecomment=[s]{/*}{*/},
  morestring=[b]',
  morestring=[b]"
}

\lstset{
  language=Java,
  basicstyle=\lstcode\footnotesize,
  keywordstyle=\color{qacodekw}\bfseries,
  ndkeywordstyle=\color{qaaccent2}\bfseries,
  identifierstyle=\color{qatext},
  stringstyle=\color{qacodestr},
  commentstyle=\color{qacodecom}\itshape,
  morecomment=[s][\color{qacodeann}]{/**}{*/},
  backgroundcolor=\color{qapanelalt},
  rulecolor=\color{qaline},
  frame=single,
  framesep=3pt,
  xleftmargin=4pt,
  xrightmargin=4pt,
  keepspaces=true,
  columns=fullflexible,
  showstringspaces=false,
  breaklines=true,
  tabsize=2,
  escapeinside={(*@}{@*)},
}

% ---------------------------------------------------------------
% CARD-STYLE BOXES
% ---------------------------------------------------------------
\tcbset{
  figcard/.style={colback=white, colframe=white, boxrule=0pt, arc=4pt,
    left=3pt, right=3pt, top=3pt, bottom=3pt, boxsep=0pt,
    halign=center, valign=center},
  datacard/.style={colback=qapanel, colframe=qaline, boxrule=0.4pt, arc=3pt,
    left=8pt, right=8pt, top=6pt, bottom=6pt, coltext=qatext,
    before upper={\color{qatext}}, halign=center},
  codebox/.style={colback=qapanelalt, colframe=qaline, boxrule=0.4pt, arc=3pt,
    left=8pt, right=8pt, top=5pt, bottom=5pt, listing only,
    listing options={basicstyle=\ttfamily\small\color{qaaccent},
      breaklines=true, showstringspaces=false, columns=fullflexible,
      keepspaces=true, frame=none, backgroundcolor={}}}
}
\newtcolorbox{figcard}[1][]{figcard,#1}
\newtcolorbox{datacard}{datacard}
\newtcblisting{codebox}{codebox}

% ---------------------------------------------------------------
% HANDY MACROS
% ---------------------------------------------------------------
\newcommand{\hl}[1]{\alert{#1}}
\newcommand{\source}[1]{{\scriptsize\color{qamuted}#1}}

% Sommaire automatique en debut de section
\AtBeginSection[]{%
  \begin{frame}<beamer>{Plan du cours}
    \tableofcontents[currentsection]
  \end{frame}}

% Page de titre : \rtitleslide{numero de chapitre}{titre du chapitre}
% Les auteurs sont repris de \rauthors (modifiable avant \begin{document})
\newcommand{\rauthors}{Selena Lamari et Chouki Tibermacine}
\newcommand{\rtitleslide}[2]{%
\begin{frame}[plain,noframenumbering]
  \begin{center}
    \vspace{2mm}
    {\color{qaaccent}\rule{2.1cm}{1.2pt}}\\[6mm]
    {\Large\bfseries\color{qatext} R3.07}\\[1mm]
    {\Large\bfseries\color{qatext} SQL dans un langage de programmation}\\[5mm]
    {\large\color{qamuted} Chapitre #1}\\[1mm]
    {\large\color{qaaccent} #2}\\[6mm]
    \includegraphics[height=2.1cm]{iut_logo_transparent.png}

    \vspace{4mm}

    {\small\color{qatext} \rauthors}\\[1mm]
    {\scriptsize\color{qamuted} BUT 2 Informatique --- IUT de Vannes --- Universit\'e Bretagne Sud}
  \end{center}
\end{frame}
}

% ---- Diapo de fin « Questions ? » (même modèle que R5.A.08) ----
\newcommand{\rcontact}{Chouki.Tibermacine@univ-ubs.fr}
\newcommand{\rquestionslide}{%
\begin{frame}{Questions ?}
  \begin{center}
    \vspace{2mm}
    \begin{tikzpicture}
      \node[regular polygon, regular polygon sides=6, draw=qaline, thick,
        minimum size=4.6cm, rotate=90] (hexin) {};
      \node[regular polygon, regular polygon sides=6, draw=qaaccent, dashed,
        thick, minimum size=6.2cm, rotate=90] (hexout) {};
      \node[font=\fontsize{92}{92}\selectfont\bfseries, color=qaaccent2] at (hexin) {?};
    \end{tikzpicture}\\[5mm]
    {\small\color{qamuted} \rcontact}
  \end{center}
\end{frame}}

\graphicspath{{../../latex/common/assets/}}

\title{Patron Repository et Spring Data JPA}
\author[C. Tibermacine]{Selena Lamari et Chouki Tibermacine}
\date{}

% SQL genere par Hibernate (sortie de spring.jpa.show-sql)
\lstdefinestyle{log}{language=SQL, basicstyle=\lstcode\scriptsize, backgroundcolor=\color{qapanel},
  keywordstyle=\color{qacodekw}, commentstyle=\color{qacodecom}}

\tikzset{
  comp/.style={rectangle, rounded corners, draw=qaaccent, thick, fill=qapanel, text=qatext,
    minimum width=3.2cm, minimum height=0.8cm, align=center, font=\scriptsize},
  itf/.style={comp, draw=qaaccent2},
  dep/.style={-Latex, thick, qamuted},
  real/.style={-{Latex[open]}, thick, dashed, qamuted}
}

\begin{document}

\rtitleslide{6}{Patron Repository et Spring Data JPA}

% ================================================================
\section{Du DAO au Repository}

\begin{frame}{O\`u en est-on ?}
  \begin{datacard}
    \small
    \begin{tabular}{@{}lp{5.2cm}p{5.2cm}@{}}
      \toprule
      \textbf{\'Etape} & \textbf{Ce qu'on \'ecrit} & \textbf{Ce qui est g\'en\'er\'e} \\
      \midrule
      JDBC (ch.\,1) & SQL, \texttt{PreparedStatement}, mapping \`a la main & --- \\
      DAO (ch.\,5) & une classe par table, SQL ou JPQL & --- \\
      JPA (ch.\,4) & entit\'es annot\'ees, JPQL & le SQL CRUD \\
      Repository (TP5) & une interface \textbf{m\'etier} + son impl\'ementation JPA & le SQL CRUD \\
      \alert{Spring Data} & \alert{l'interface seulement} & le SQL \textbf{et} l'impl\'ementation \\
      \bottomrule
    \end{tabular}
  \end{datacard}
  \vspace{2mm}
  \begin{itemize}
    \item DAO : vocabulaire \textbf{technique}, une classe par table (\texttt{insertClient}, \texttt{updateClient})
    \item Repository : une \textbf{collection d'objets m\'etier} dont on ignore le stockage (\texttt{save}, \texttt{findById}, \texttt{findComptesDuClient})
  \end{itemize}
\end{frame}

\begin{frame}{Le Repository inverse la d\'ependance}
  \begin{center}
  \begin{tikzpicture}
    \node[comp] (svc) at (0,0) {BanqueService\\{\tiny(m\'etier)}};
    \node[itf]  (rep) at (5,0) {\guillemotleft interface\guillemotright\\ClientRepository};
    \node[comp] (jpa) at (2.6,-2) {ClientRepositoryJPA\\{\tiny(EntityManager)}};
    \node[comp] (mem) at (7.4,-2) {ClientRepositoryMemoire\\{\tiny(HashMap, pour les tests)}};
    \draw[dep] (svc) -- node[above, font=\tiny, text=qamuted] {utilise} (rep);
    \draw[real] (jpa) -- (rep);
    \draw[real] (mem) -- (rep);
  \end{tikzpicture}
  \end{center}
  \begin{itemize}
    \item Le service ne conna\^it que l'\textbf{interface} : ni SQL, ni JPA, ni \texttt{EntityManager}
    \item On peut changer de technique de persistance sans toucher au m\'etier
    \item On peut \textbf{tester le m\'etier sans base de donn\'ees}
  \end{itemize}
\end{frame}

\begin{frame}[fragile]{Tester le m\'etier sans base}
\begin{lstlisting}[language=Java]
class ClientRepositoryMemoire implements ClientRepository {
  private final Map<Integer, Client> donnees = new HashMap<>();
  public void save(Client c)                 { donnees.put(c.getId(), c); }
  public Optional<Client> findById(int id)   { return Optional.ofNullable(donnees.get(id)); }
  public List<Client> findAll()              { return new ArrayList<>(donnees.values()); }
  public List<Compte> findComptesByClient(int id) {
    return findById(id).map(Client::getComptes).orElse(List.of());
  }
}

@Test void soldeTotalDUnClient() {
  ClientRepository repo = new ClientRepositoryMemoire();
  repo.save(clientAvecDeuxComptes(1, 100, 50));
  assertEquals(150, new BanqueService(repo).soldeTotal(1));   // aucun Docker, aucun SQL
}
\end{lstlisting}
\end{frame}

\begin{frame}[fragile]{Ce qui se r\'ep\`ete d'un Repository \`a l'autre}
  \begin{columns}[T]
    \begin{column}{0.5\textwidth}
\begin{lstlisting}[language=Java]
public Optional<Client> findById(int id) {
  return Optional.ofNullable(
      em.find(Client.class, id));
}
public List<Client> findByVille(String v) {
  return em.createQuery(
      "SELECT c FROM Client c WHERE c.ville = :v",
      Client.class)
    .setParameter("v", v).getResultList();
}
\end{lstlisting}
    \end{column}
    \begin{column}{0.46\textwidth}
      \begin{itemize}
        \item Le m\^eme code pour \texttt{Client}, \texttt{Compte}, \texttt{Operation}\ldots
        \item La requ\^ete se d\'eduit du \textbf{nom} de la m\'ethode
        \item[$\Rightarrow$] Un framework peut l'\'ecrire \`a notre place : \alert{Spring Data}
      \end{itemize}
      \vspace{2mm}
      {\small\color{qamuted} Dans le TP5, vous \'ecrivez cette impl\'ementation \`a la main pour comprendre ce que Spring Data g\'en\`ere.}
    \end{column}
  \end{columns}
\end{frame}

% ================================================================
\section{Spring Data JPA}

\begin{frame}[fragile]{Mise en place (Spring Boot)}
  \begin{columns}[T]
    \begin{column}{0.5\textwidth}
      \textbf{\texttt{pom.xml}}
\begin{lstlisting}[language=XML]
<dependency>
  <groupId>org.springframework.boot</groupId>
  <artifactId>spring-boot-starter-data-jpa</artifactId>
</dependency>
<dependency>
  <groupId>org.postgresql</groupId>
  <artifactId>postgresql</artifactId>
  <scope>runtime</scope>
</dependency>
\end{lstlisting}
    \end{column}
    \begin{column}{0.48\textwidth}
      \textbf{\texttt{application.properties}}
\begin{lstlisting}[language={}]
spring.datasource.url=jdbc:postgresql://localhost:5432/banque_db
spring.datasource.username=${DB_USER}
spring.datasource.password=${DB_PASSWORD}
spring.jpa.hibernate.ddl-auto=validate
spring.jpa.show-sql=true
\end{lstlisting}
      \begin{itemize}
        \item Plus de \texttt{persistence.xml} ni d'\texttt{EntityManagerFactory} \`a cr\'eer
        \item Pool de connexions (HikariCP) fourni
        \item Les entit\'es JPA du chapitre 4 sont r\'eutilis\'ees telles quelles
      \end{itemize}
    \end{column}
  \end{columns}
\end{frame}

\begin{frame}[fragile]{D\'eclarer un Repository : une interface suffit}
\begin{lstlisting}[language=Java]
public interface ClientRepository extends JpaRepository<Client, Integer> {
  // rien a ecrire pour le CRUD
}

@Service
public class BanqueService {
  private final ClientRepository clients;
  public BanqueService(ClientRepository clients) {   // injection par constructeur
    this.clients = clients;
  }
}
\end{lstlisting}
  \begin{itemize}
    \item Au d\'emarrage, Spring \textbf{g\'en\`ere une classe} qui impl\'emente l'interface et l'injecte
    \item M\'ethodes h\'erit\'ees : \texttt{save}, \texttt{saveAll}, \texttt{findById}, \texttt{findAll}, \texttt{existsById}, \texttt{count},
      \texttt{deleteById}, \texttt{findAll(Pageable)}, \texttt{findAll(Sort)}\ldots
    \item \texttt{Client} et \texttt{Integer} : type de l'entit\'e et type de sa cl\'e primaire
  \end{itemize}
\end{frame}

\begin{frame}[fragile]{Requ\^etes d\'eriv\'ees du nom de la m\'ethode}
\begin{lstlisting}[language=Java]
List<Client>  findByVille(String ville);
List<Client>  findByNomClientStartingWithOrderByNomClientAsc(String prefixe);
long          countByVille(String ville);
boolean       existsByNomClient(String nom);
List<Compte>  findByClientNumClient(int numClient);         // dans CompteRepository
List<Compte>  findBySoldeGreaterThanEqual(BigDecimal min);
\end{lstlisting}
  SQL r\'eellement g\'en\'er\'e par Hibernate (extraits du journal) :
\begin{lstlisting}[style=log]
select ... from client c1_0 where c1_0.ville=?
select ... from client c1_0 where c1_0.nom_client like ? escape '\' order by c1_0.nom_client
select count(c1_0.num_client) from client c1_0 where c1_0.ville=?
select ... from compte c1_0 left join client c2_0 on c2_0.num_client=c1_0.num_client where c2_0.num_client=?
\end{lstlisting}
  \begin{itemize}
    \item Mots-cl\'es : \texttt{And}, \texttt{Or}, \texttt{Between}, \texttt{LessThan}, \texttt{Like}, \texttt{In}, \texttt{IsNull}, \texttt{OrderBy}, \texttt{First10}\ldots
    \item Un nom mal form\'e ou une propri\'et\'e inconnue : erreur \textbf{au d\'emarrage}, pas \`a l'ex\'ecution
  \end{itemize}
\end{frame}

\begin{frame}[fragile]{Quand le nom devient trop long : \texttt{@Query}}
\begin{lstlisting}[language=Java]
@Query("SELECT c FROM Client c WHERE size(c.comptes) >= :n")          // JPQL
List<Client> avecAuMoins(@Param("n") int n);

@Query(value = "SELECT * FROM client WHERE ville = ?1 LIMIT 5",      // SQL natif
       nativeQuery = true)
List<Client> cinqClientsDe(String ville);

@Modifying                                                           // UPDATE / DELETE
@Query("UPDATE Compte c SET c.solde = c.solde - :m WHERE c.numCompte = :num AND c.solde >= :m")
int debiter(@Param("num") int num, @Param("m") BigDecimal m);         // nb de lignes modifiees
\end{lstlisting}
  \begin{itemize}
    \item \texttt{@Modifying} : mise \`a jour relative du chapitre 2 (\texttt{solde = solde - ?}), ex\'ecut\'ee dans une transaction
    \item \alert{Attention} : une requ\^ete \texttt{@Modifying} \textbf{contourne} le contexte de persistance et le compteur \texttt{@Version}
      (v\'erifi\'e : la version du compte reste \`a 0)
  \end{itemize}
\end{frame}

\begin{frame}[fragile]{Projections : ne charger que l'utile (DTO, chapitre 5)}
\begin{lstlisting}[language=Java]
public interface ClientResume {          // projection : seulement des getters
  String getNomClient();
  String getVille();
}
List<ClientResume> findByVilleOrderByNomClient(String ville);
\end{lstlisting}
\begin{lstlisting}[style=log]
select c1_0.nom_client,c1_0.ville from client c1_0 where c1_0.ville=? order by 1
\end{lstlisting}
  \begin{itemize}
    \item Seules les colonnes demand\'ees sont lues (chapitre 3 : \og pas de \texttt{SELECT *} \fg{})
    \item Variante : un \texttt{record} Java comme type de retour
    \item Le service renvoie ces projections ou des DTO, jamais les entit\'es, \`a la couche pr\'esentation
  \end{itemize}
\end{frame}

\begin{frame}[fragile]{Pagination : \texttt{Page} (OFFSET) ou \texttt{Window} (par cl\'e)}
  \textbf{Par num\'ero de page}
\begin{lstlisting}[language=Java,basicstyle=\lstcode\scriptsize]
Page<Client> findByVille(String v, Pageable p);
Page<Client> page = repo.findByVille("Vannes", PageRequest.of(1, 20, Sort.by("nomClient")));
\end{lstlisting}
\begin{lstlisting}[style=log]
select ... where c1_0.ville=? order by c1_0.nom_client offset ? rows fetch first ? rows only
\end{lstlisting}
  \textbf{Par cl\'e (Spring Data 3.1+)}
\begin{lstlisting}[language=Java,basicstyle=\lstcode\scriptsize]
Window<Client> findFirst20ByVilleOrderByNumClient(String v, ScrollPosition pos);
Window<Client> w = repo.findFirst20ByVilleOrderByNumClient("Vannes", ScrollPosition.keyset());
w = repo.findFirst20ByVilleOrderByNumClient("Vannes", w.positionAt(w.size() - 1));  // page suivante
\end{lstlisting}
\begin{lstlisting}[style=log]
select ... where c1_0.ville=? and (c1_0.num_client>?) order by c1_0.num_client fetch first ? rows only
\end{lstlisting}
  {\small \texttt{Page} est pratique (num\'ero de page, nombre total) mais repose sur \texttt{OFFSET}, lent sur les pages lointaines
  (chapitre 3). \texttt{Window} produit la pagination par cl\'e.}
\end{frame}

% ================================================================
\section{Transactions et concurrence avec Spring}

\begin{frame}[fragile]{\texttt{@Transactional} sur la couche service}
\begin{lstlisting}[language=Java]
@Service
public class BanqueService {
  @Transactional                                        // BEGIN ... COMMIT / ROLLBACK
  public void virer(int src, int dst, BigDecimal montant) {
    if (comptes.debiter(src, montant) == 0) throw new SoldeInsuffisantException(src);
    comptes.crediter(dst, montant);
    operations.save(new Operation(src, dst, montant));
  }
  @Transactional(readOnly = true)
  public BigDecimal soldeTotal(int numClient) { ... }
}
\end{lstlisting}
  \begin{itemize}
    \item Remplace le sch\'ema \texttt{setAutoCommit(false)} / \texttt{commit()} / \texttt{rollback()} du chapitre 2
    \item Rollback par d\'efaut sur les exceptions \textbf{non v\'erifi\'ees} (\texttt{RuntimeException}) seulement ;
      sinon : \texttt{@Transactional(rollbackFor = Exception.class)}
    \item Niveau d'isolation : \texttt{@Transactional(isolation = Isolation.REPEATABLE\_READ)}
    \item Sans effet sur une m\'ethode \texttt{private} et lors d'un appel interne (\texttt{this.virer(\ldots)}) : la transaction est pos\'ee par un \textbf{proxy} autour du bean
  \end{itemize}
\end{frame}

\begin{frame}[fragile]{Verrou optimiste : \texttt{@Version}}
  \begin{columns}[T]
    \begin{column}{0.42\textwidth}
\begin{lstlisting}[language=Java]
@Entity
public class Compte {
  @Id private Integer numCompte;
  private BigDecimal solde;
  @Version private int version;
  ...
}
\end{lstlisting}
    \end{column}
    \begin{column}{0.55\textwidth}
\begin{lstlisting}[style=log]
update compte set num_client=?,solde=?,version=?
 where num_compte=? and version=?
\end{lstlisting}
      Si la ligne a chang\'e depuis la lecture :\\[1mm]
      \texttt{\footnotesize ObjectOptimisticLockingFailureException:}\\
      \texttt{\footnotesize Row was updated or deleted by another transaction}
    \end{column}
  \end{columns}
  \vspace{2mm}
  \begin{itemize}
    \item C'est la \og solution 3 \fg{} du chapitre 2, \'ecrite par Hibernate
    \item \`A attraper dans le contr\^oleur : message \og donn\'ees modifi\'ees entre-temps, rechargez \fg{} ou nouvelle tentative
  \end{itemize}
\end{frame}

\begin{frame}[fragile]{Verrou pessimiste : \texttt{@Lock}}
\begin{lstlisting}[language=Java]
@Lock(LockModeType.PESSIMISTIC_WRITE)
@Query("SELECT c FROM Compte c WHERE c.numCompte = :num")
Optional<Compte> findPourMaj(@Param("num") int num);
\end{lstlisting}
\begin{lstlisting}[style=log]
select ... from compte c1_0 where c1_0.num_compte=? for no key update
\end{lstlisting}
  \begin{itemize}
    \item \og Solution 2 \fg{} du chapitre 2. Sur PostgreSQL, Hibernate produit \texttt{FOR NO KEY UPDATE}, une variante de
      \texttt{FOR UPDATE} qui ne bloque pas les insertions de lignes filles r\'ef\'eren\c{c}ant ce compte
    \item Appel\'ee hors transaction : \texttt{InvalidDataAccessApiUsageException: Query requires transaction be in progress}
  \end{itemize}
\end{frame}

% ================================================================
\section{Performances}

\begin{frame}[fragile]{Le probl\`eme N+1 et \texttt{@EntityGraph}}
  \begin{columns}[T]
    \begin{column}{0.49\textwidth}
\begin{lstlisting}[language=Java]
for (Client c : repo.findByVille("Vannes"))
  System.out.println(c.getComptes().size());
\end{lstlisting}
\begin{lstlisting}[style=log]
select ... from client c1_0 where c1_0.ville=?
select ... from compte c1_0 where c1_0.num_client=?
select ... from compte c1_0 where c1_0.num_client=?
select ... from compte c1_0 where c1_0.num_client=?
\end{lstlisting}
      \centering\alert{1 + N requ\^etes}
    \end{column}
    \begin{column}{0.49\textwidth}
\begin{lstlisting}[language=Java]
@EntityGraph(attributePaths = "comptes")
List<Client> findWithComptesByVille(String v);
\end{lstlisting}
\begin{lstlisting}[style=log]
select ... from client c1_0
  left join compte c2_0
    on c1_0.num_client=c2_0.num_client
 where c1_0.ville=?
\end{lstlisting}
      \centering\textbf{1 requ\^ete}
    \end{column}
  \end{columns}
  \vspace{2mm}
  {\small Toujours activer \texttt{spring.jpa.show-sql} en d\'eveloppement et \textbf{compter} les requ\^etes. Alternative : \texttt{JOIN FETCH} dans un \texttt{@Query}.}
\end{frame}

\begin{frame}[fragile]{Quand revenir au SQL}
  \begin{itemize}
    \item Requ\^etes analytiques complexes (fen\^etres, CTE r\'ecursives) : \texttt{@Query(nativeQuery = true)}
    \item Traitements de masse : un seul \texttt{UPDATE} plut\^ot que charger 100\,000 entit\'es
    \item Logique ex\'ecut\'ee dans la base (chapitre 7) :
  \end{itemize}
\begin{lstlisting}[language=Java]
@Modifying
@Query(value = "CALL virement(:src, :dst, CAST(:m AS NUMERIC))", nativeQuery = true)
void virement(@Param("src") int src, @Param("dst") int dst, @Param("m") BigDecimal m);
\end{lstlisting}
  \begin{itemize}
    \item Spring Data ne masque pas la base : il faut savoir lire le SQL qu'il produit (chapitres 1 \`a 3)
  \end{itemize}
\end{frame}

\begin{frame}{\`A retenir}
  \begin{itemize}
    \item Repository = interface \textbf{m\'etier} ; le service en d\'epend, jamais de JPA
    \item Spring Data g\'en\`ere l'impl\'ementation : CRUD h\'erit\'e, requ\^etes d\'eriv\'ees, \texttt{@Query}
    \item Projections pour ne lire que l'utile ; \texttt{Window} pour paginer sans \texttt{OFFSET}
    \item \texttt{@Transactional} sur le service, \texttt{@Version} (optimiste) ou \texttt{@Lock} (pessimiste) contre les mises \`a jour perdues
    \item \texttt{@EntityGraph} ou \texttt{JOIN FETCH} contre le N+1
    \item Toujours v\'erifier le SQL g\'en\'er\'e
  \end{itemize}
\end{frame}

\rquestionslide

\end{document}
